% GENERATED FILE. Edit canonical lessons, then run npm run generate:books.
% canonical-source-hash: fnv1a64:04fb2542ff522c68
% canonical-lessons: KA-C252-adrsta, KA-C252-salahe, KA-C252-sumaru, KA-C252-jotege, KA-C252-ottige, KA-R252-first-pass-words-from-chapters-244-247, KA-R252-second-pass-words-from-chapters-248-252

\chapter{Luck, Advice, Approximately, Along with, Together}
\label{ch:ka-luck-advice-approximately-along-with-together}
\hlchaptermodality{252}

\begin{chapteropening}
\textbf{By the end of this chapter:} \emph{I can say luck, advice, approximately, along with and together.}

Complete the last lesson of chapter 252: I can say luck, advice, approximately, along with and together.
\end{chapteropening}

\section[\texorpdfstring{\kn{ಅದೃಷ್ಟ} (adṛṣṭa)}{ಅದೃಷ್ಟ (adṛṣṭa)}]{\kn{ಅದೃಷ್ಟ} (adṛṣṭa) — luck}

\label{lesson:KA-C252-adrsta}



\begin{quote}
Before the new one: say the Kannada for a downward slope, then the Kannada for a backyard.
\end{quote}

\subsection*{You'll want to know: \kn{ಅದೃಷ್ಟ}}
\textbf{\kn{ಅದೃಷ್ಟ}} — \emph{adṛṣṭa} — \textquotedblleft{}luck\textquotedblright{}.

\begin{grammarlens}[title={What you've built}]
One more word for this chapter.
\end{grammarlens}

\subsection*{Your turn}
\begin{itemize}
\raggedright
  \item \emph{Say it:} \emph{adṛṣṭa}
  \item \emph{Say it:} \emph{adṛṣṭa}, once more
  \item \emph{Say it:} say \emph{adṛṣṭa}
  \item \emph{Recall:} say the Kannada for a downward slope, then the Kannada for a backyard, then say \emph{adṛṣṭa} again
\end{itemize}

\begin{tcolorbox}[breakable,colback=teal!4,colframe=teal!35!black,title={Before you move on}]
What does \kn{ಅದೃಷ್ಟ} mean? (\textquotedblleft{}Luck\textquotedblright{}.) Say it once more.
\end{tcolorbox}
\section[\texorpdfstring{\kn{ಸಲಹೆ} (salahe)}{ಸಲಹೆ (salahe)}]{\kn{ಸಲಹೆ} (salahe) — advice, a suggestion}

\label{lesson:KA-C252-salahe}



\begin{quote}
Before the new one: say the Kannada for a backyard, then the Kannada for luck.
\end{quote}

\subsection*{You'll want to know: \kn{ಸಲಹೆ}}
\textbf{\kn{ಸಲಹೆ}} — \emph{salahe} — \textquotedblleft{}advice, a suggestion\textquotedblright{}.

\begin{grammarlens}[title={What you've built}]
One more word for this chapter.
\end{grammarlens}

\subsection*{Your turn}
\begin{itemize}
\raggedright
  \item \emph{Say it:} \emph{salahe}
  \item \emph{Say it:} \emph{salahe}, once more
  \item \emph{Say it:} say \emph{salahe}
  \item \emph{Recall:} say the Kannada for a backyard, then the Kannada for luck, then say \emph{salahe} again
\end{itemize}

\begin{tcolorbox}[breakable,colback=teal!4,colframe=teal!35!black,title={Before you move on}]
What does \kn{ಸಲಹೆ} mean? (\textquotedblleft{}Advice, a suggestion\textquotedblright{}.) Say it once more.
\end{tcolorbox}
\section[\texorpdfstring{\kn{ಸುಮಾರು} (sumāru)}{ಸುಮಾರು (sumāru)}]{\kn{ಸುಮಾರು} (sumāru) — approximately, roughly}

\label{lesson:KA-C252-sumaru}



\begin{quote}
Before the new one: say the Kannada for luck, then the Kannada for advice.
\end{quote}

\subsection*{You'll want to know: \kn{ಸುಮಾರು}}
\textbf{\kn{ಸುಮಾರು}} — \emph{sumāru} — \textquotedblleft{}approximately, roughly\textquotedblright{}.

\begin{grammarlens}[title={What you've built}]
One more word for this chapter.
\end{grammarlens}

\subsection*{Your turn}
\begin{itemize}
\raggedright
  \item \emph{Say it:} \emph{sumāru}
  \item \emph{Say it:} \emph{sumāru}, once more
  \item \emph{Say it:} say \emph{sumāru}
  \item \emph{Recall:} say the Kannada for luck, then the Kannada for advice, then say \emph{sumāru} again
\end{itemize}

\begin{tcolorbox}[breakable,colback=teal!4,colframe=teal!35!black,title={Before you move on}]
What does \kn{ಸುಮಾರು} mean? (\textquotedblleft{}Approximately, roughly\textquotedblright{}.) Say it once more.
\end{tcolorbox}
\section[\texorpdfstring{\kn{ಜೊತೆಗೆ} (jotege)}{ಜೊತೆಗೆ (jotege)}]{\kn{ಜೊತೆಗೆ} (jotege) — along with}

\label{lesson:KA-C252-jotege}



\begin{quote}
Before the new one: say the Kannada for advice, then the Kannada for approximately.
\end{quote}

\subsection*{You'll want to know: \kn{ಜೊತೆಗೆ}}
\textbf{\kn{ಜೊತೆಗೆ}} — \emph{jotege} — \textquotedblleft{}along with\textquotedblright{}.

\begin{grammarlens}[title={What you've built}]
One more word for this chapter.
\end{grammarlens}

\subsection*{Your turn}
\begin{itemize}
\raggedright
  \item \emph{Say it:} \emph{jotege}
  \item \emph{Say it:} \emph{jotege}, once more
  \item \emph{Say it:} say \emph{jotege}
  \item \emph{Recall:} say the Kannada for advice, then the Kannada for approximately, then say \emph{jotege} again
\end{itemize}

\begin{tcolorbox}[breakable,colback=teal!4,colframe=teal!35!black,title={Before you move on}]
What does \kn{ಜೊತೆಗೆ} mean? (\textquotedblleft{}Along with\textquotedblright{}.) Say it once more.
\end{tcolorbox}
\section[\texorpdfstring{\kn{ಒಟ್ಟಿಗೆ} (oṭṭige)}{ಒಟ್ಟಿಗೆ (oṭṭige)}]{\kn{ಒಟ್ಟಿಗೆ} (oṭṭige) — together}

\label{lesson:KA-C252-ottige}



\begin{quote}
Before the new one: say the Kannada for approximately, then the Kannada for along with.
\end{quote}

\subsection*{You'll want to know: \kn{ಒಟ್ಟಿಗೆ}}
\textbf{\kn{ಒಟ್ಟಿಗೆ}} — \emph{oṭṭige} — \textquotedblleft{}together\textquotedblright{}.

\begin{grammarlens}[title={What you've built}]
One more word for this chapter.
\end{grammarlens}

\subsection*{Your turn}
\begin{itemize}
\raggedright
  \item \emph{Say it:} \emph{oṭṭige}
  \item \emph{Say it:} \emph{oṭṭige}, once more
  \item \emph{Say it:} say \emph{oṭṭige}
  \item \emph{Recall:} say the Kannada for approximately, then the Kannada for along with, then say \emph{oṭṭige} again
\end{itemize}

\begin{tcolorbox}[breakable,colback=teal!4,colframe=teal!35!black,title={Before you move on}]
What does \kn{ಒಟ್ಟಿಗೆ} mean? (\textquotedblleft{}Together\textquotedblright{}.) Say it once more.
\end{tcolorbox}
\section[(dialogue)]{Words from chapters 244-247 — a first pass}

\label{lesson:KA-R252-first-pass-words-from-chapters-244-247}



\begin{quote}
Say the words for along with and together.
\end{quote}

\subsection*{Your turn}
Say these aloud:

\begin{itemize}
\raggedright
  \item to move aside, to shrink, to lick, to suck, to install
  \item diarrhoea, an allergy, a swelling, a pimple, medical treatment
  \item a sprain, indigestion, a sore, a bandage, diabetes
  \item Dasara, Deepavali, Sankranti, Ganesha Chaturthi, Rajyotsava
\end{itemize}

\begin{tcolorbox}[breakable,colback=teal!4,colframe=teal!35!black,title={Before you move on}]
To move aside? (\textbf{\kn{ಸರಿಸು}}, \emph{sarisu}.) To shrink? (\textbf{\kn{ಕುಗ್ಗು}}, \emph{kuggu}.) To lick? (\textbf{\kn{ನೆಕ್ಕು}}, \emph{nekku}.) To suck? (\textbf{\kn{ಹೀರು}}, \emph{hīru}.) To install? (\textbf{\kn{ಅಳವಡಿಸು}}, \emph{aḷavaḍisu}.) Diarrhoea? (\textbf{\kn{ಭೇದಿ}}, \emph{bhēdi}.) An allergy? (\textbf{\kn{ಅಲರ್ಜಿ}}, \emph{alarji}.) A swelling? (\textbf{\kn{ಊತ}}, \emph{ūta}.) A pimple? (\textbf{\kn{ಗುಳ್ಳೆ}}, \emph{guḷḷe}.) Medical treatment? (\textbf{\kn{ಚಿಕಿತ್ಸೆ}}, \emph{cikitse}.) A sprain? (\textbf{\kn{ಉಳುಕು}}, \emph{uḷuku}.) Indigestion? (\textbf{\kn{ಅಜೀರ್ಣ}}, \emph{ajīrṇa}.) A sore? (\textbf{\kn{ಹುಣ್ಣು}}, \emph{huṇṇu}.) A bandage? (\textbf{\kn{ಬ್ಯಾಂಡೇಜ್}}, \emph{byāṇḍēj}.) Diabetes? (\textbf{\kn{ಮಧುಮೇಹ}}, \emph{madhumēha}.) Dasara? (\textbf{\kn{ದಸರಾ}}, \emph{dasarā}.) Deepavali? (\textbf{\kn{ದೀಪಾವಳಿ}}, \emph{dīpāvaḷi}.) Sankranti? (\textbf{\kn{ಸಂಕ್ರಾಂತಿ}}, \emph{saṅkrānti}.) Ganesha Chaturthi? (\textbf{\kn{ಗಣೇಶ} \kn{ಚತುರ್ಥಿ}}, \emph{gaṇēśa caturthi}.) Rajyotsava? (\textbf{\kn{ರಾಜ್ಯೋತ್ಸವ}}, \emph{rājyōtsava}.)
\end{tcolorbox}
\section[(dialogue)]{Words from chapters 248-252 — a second pass}

\label{lesson:KA-R252-second-pass-words-from-chapters-248-252}



\begin{quote}
Say the words for a decade and an anniversary.
\end{quote}

\subsection*{Your turn}
Say these aloud:

\begin{itemize}
\raggedright
  \item a decade, an anniversary, a deadline, an auspicious time, one and a half
  \item east, west, south, opposite, a district
  \item around, a hut, a residential area, a post office, a valley
  \item a frontier, a canal, a tunnel, a downward slope, a backyard
  \item luck, advice, approximately, along with, together
\end{itemize}

\begin{tcolorbox}[breakable,colback=teal!4,colframe=teal!35!black,title={Before you move on}]
A decade? (\textbf{\kn{ದಶಕ}}, \emph{daśaka}.) An anniversary? (\textbf{\kn{ವಾರ್ಷಿಕೋತ್ಸವ}}, \emph{vārṣikōtsava}.) A deadline? (\textbf{\kn{ಗಡುವು}}, \emph{gaḍuvu}.) An auspicious time? (\textbf{\kn{ಮುಹೂರ್ತ}}, \emph{muhūrta}.) One and a half? (\textbf{\kn{ಒಂದೂವರೆ}}, \emph{ondūvare}.) East? (\textbf{\kn{ಪೂರ್ವ}}, \emph{pūrva}.) West? (\textbf{\kn{ಪಶ್ಚಿಮ}}, \emph{paścima}.) South? (\textbf{\kn{ದಕ್ಷಿಣ}}, \emph{dakṣiṇa}.) Opposite? (\textbf{\kn{ಎದುರು}}, \emph{eduru}.) A district? (\textbf{\kn{ಜಿಲ್ಲೆ}}, \emph{jille}.) Around? (\textbf{\kn{ಸುತ್ತ}}, \emph{sutta}.) A hut? (\textbf{\kn{ಗುಡಿಸಲು}}, \emph{guḍisalu}.) A residential area? (\textbf{\kn{ಬಡಾವಣೆ}}, \emph{baḍāvaṇe}.) A post office? (\textbf{\kn{ಅಂಚೆ} \kn{ಕಚೇರಿ}}, \emph{añce kacēri}.) A valley? (\textbf{\kn{ಕಣಿವೆ}}, \emph{kaṇive}.) A frontier? (\textbf{\kn{ಗಡಿ}}, \emph{gaḍi}.) A canal? (\textbf{\kn{ಕಾಲುವೆ}}, \emph{kāluve}.) A tunnel? (\textbf{\kn{ಸುರಂಗ}}, \emph{suraṅga}.) A downward slope? (\textbf{\kn{ಇಳಿಜಾರು}}, \emph{iḷijāru}.) A backyard? (\textbf{\kn{ಹಿತ್ತಿಲು}}, \emph{hittilu}.) Luck? (\textbf{\kn{ಅದೃಷ್ಟ}}, \emph{adṛṣṭa}.) Advice? (\textbf{\kn{ಸಲಹೆ}}, \emph{salahe}.) Approximately? (\textbf{\kn{ಸುಮಾರು}}, \emph{sumāru}.) Along with? (\textbf{\kn{ಜೊತೆಗೆ}}, \emph{jotege}.) Together? (\textbf{\kn{ಒಟ್ಟಿಗೆ}}, \emph{oṭṭige}.)
\end{tcolorbox}
